%
% Copyright assignment to GNUnet e.V. in der Fassung vom 20.1.2016
% Inhaltliche Korrekturen von ng0 bis zum 14.03 2019:
% - URL changes (git repository location)
% - Lizenzaenderung (GPL3 -> AGPL3)
%
\documentclass[a4paper]{article} %,titlepage
%\documentclass{article}
\usepackage{german,url,a4wide}
\usepackage[utf8]{inputenc}
\usepackage[T1]{fontenc}
\usepackage{graphicx}
\usepackage[%
  pdftitle = {Copyright assignment},
%%  pdfsubject = {},
  pdfpagemode=UseOutlines,
  colorlinks=false,
%%    linkcolor=black,
%%    filecolor=black,
%%    urlcolor=black,
%%    citecolor=black
  pdftex=true,
  plainpages=false,
  hypertexnames=true,%false,
  pdfpagelabels=true,
  hyperindex=true%
]{hyperref}


\newcommand{\zweidrittel}{2/3}
\renewcommand{\thesection}{\S\arabic{section}}
\renewcommand{\thepart}{\alph{part})}
%begin{latexonly}
\makeatletter
\renewcommand\l@section[2]{%
  \ifnum \c@tocdepth >\z@
    \addpenalty\@secpenalty
    \addvspace{1.0em \@plus\p@}%
    \setlength\@tempdima{2.0em}%
    \begingroup
      \parindent \z@ \rightskip \@pnumwidth
      \parfillskip -\@pnumwidth
      \leavevmode \bfseries
      \advance\leftskip\@tempdima
      \hskip -\leftskip
      #1\nobreak\hfil \nobreak\hb@xt@\@pnumwidth{\hss #2}\par
    \endgroup
  \fi}
\makeatother
%end{latexonly}

\newenvironment{enumpar}{%
%begin{latexonly}
  \renewcommand{\labelenumi}{\theenumi}%
  \renewcommand{\theenumi}{(\arabic{enumi})}%
  \renewcommand{\labelenumii}{\theenumii}%
  \renewcommand{\theenumii}{(\alph{enumii})}%
%end{latexonly}
  \begin{enumerate}}
{\end{enumerate}}

\newcommand{\parref}[2]{\ref{#1}\ref{#2}}


\title{Copyright Assignment Agreement \\ for Contributors to GNUnet}
\date{}
\begin{document}
\maketitle
%\textit{\small{Aus Gründen der Lesbarkeit werden im
%    Folgenden weitgehend männliche Bezeichnungen verwendet. Natürlich
%    sind in jedem Falle weibliche wie männliche Personen gleichermaßen
%    gemeint und angesprochen. Wir bitten darum die jeweils weiblichen
%    Fassungen mitzudenken.}}
\noindent
{\bf Between:}

GNUnet e.V., a German association of GNUnet developers
registered in M\"unchen and seated in Garching bei M\"unchen,
represented by

\vspace{2cm}

an authorized currently elected representative of the
GNUnet e.V. Vorstand, hereinafter referred to as {\bf ``GNUnet e.V.''},

\noindent
{\bf And:}
% Note: pseudonymous contributions are still OK!
% Then we just put the pseudonym.

\vspace{4cm}
hereinafter referred to as {\bf ``the Contributor''}.

\vspace{0.5cm}

\noindent
{\bf IT HAVING BEEN PREVIOUSLY ESTABLISHED:}
\begin{itemize}
\item That the GNUnet project is directed and supported by GNUnet e.V.,
  which aims to develop the GNUnet software as part of the GNU project,
  distributed under license GNU AGPL v3+ on the date of signing of this
  contract;
\item That this development project is opened to contributions
  submitted by individuals outside of GNUnet e.V., contributions which
  may be officially submitted by their holder to GNUnet e.V. for the
  purpose of being integrated by the latter into the GNUnet software,
  in successive versions edited and distributed by GNUnet e.V. or the
  GNU project;
\item That GNUnet e.V. wishes, in this context, to centralize copyright
  ownership on any new contribution integrated into the GNUnet software
  edited and distributed by it;
\item That the Contributor wishes to participate, or authorize its
  personnel to participate, in the development of the aforementioned
  software.
\end{itemize}

\noindent
{\bf IT HAS BEEN AGREED AS FOLLOWS:}

\section{Definition}

\begin{itemize}
  \item {\bf ``Software'':} means the GNUnet software available at
    {\tt https://git.gnunet.org/} and on GNU's FTP mirrors, in its
    present version and future versions, distributed under the GNU AGPL
    v3 license at the date of the signing of this contract, or any
    other license chosen by GNUnet e.V.
\end{itemize}

However, ``Software'' does not mean the versions derived from the
Software, developed and distributed by the Contributor, independently
of the successive versions edited and distributed by GNUnet e.V., in
accordance with the rights granted to the Contributor by the GNU AGPL
v3+ license.

\begin{itemize}
  \item {\bf ``Contribution'':} means any original contribution
    protected by copyright, in particular modifications or the
    development of new software components and new functionalities by
    the Contributor (be the Contributor an individual or an
    organization employing individuals), and that the latter
    intentionally submits to GNUnet e.V. to be integrated into the
    Software, with GNUnet e.V.'s approval. A Contribution includes its
    source code, its object code, as well as any specifications and
    documentation related thereto.
\end{itemize}

\section{Subject}

This contract defines the conditions under which the Contributor
assigns, free of charge, the copyrights over its Contributions as they
are in progress, to GNUnet e.V..  Any participation by the Contributor
or employees of the Contributor will be made under the conditions of
this contract.

Should the Contributor wish to terminate further participation under
the terms of this contract, it will inform GNUnet e.V. of this in
advance and by registered mail to:

\begin{verbatim}
GNUnet e.V.
Fakultät für Informatik -- I8
Technische Universität München
Boltzmannstraße 3
85748 Garching
GERMANY
\end{verbatim}

It is specified that this copyright assignment is a necessary
condition to enable GNUnet e.V. to integrate the Contributions of the
Contributor’s personnel in the Software.

\section{EFFECT AND DURATION}

The contract takes effect on the latest date of signing.

The contract shall be enforced for a duration of five years, renewable
by written agreement, as far as the Contributor’s participation in the
development of the Software is concerned.

However, the Contributor’s copyright on the Contributions is assigned
to GNUnet e.V. for as long as said copyrights last.

\section{COPYRIGHTS ASSIGNED TO GNUnet e.V.}

\begin{enumpar}
  \item is stipulated that copyrights over the Contributions, as
    assigned by the Contributor to GNUnet e.V., are the following:
  \begin{itemize}
    \item The right to reproduce, permanently or temporarily, all or
      part of the Contributions, in any format, on any medium and by
      any technical means, present or future, necessary for their use.
    \item The right to represent all or part of the Contributions, in
      any place, public or not, through any form of communication,
      present or future, by the means referred to above.
    \item The right to adapt for purposes of freely carrying out, on
      any medium and by any technical means, any adaptation and/or
      modification of all or part of the Contributions, including with
      a view to developing derivative software and integrating the
      Contributions into a future version of the Software. It is
      stipulated that the right of adaptation includes the corrective
      and developmental maintenance of the software into which a
      Contribution might be integrated.  Therefore, the Contributor
      promises to deliver to GNUnet e.V. the source code of the
      Contributions transferred, as well as any documentation
      pertaining thereto.
    \item The right to distribute the Contributions referred to above,
      internally or publicly, remunerated or free of charge, directly
      or through a provider and/or licensee and/or assignee.
    \item The right of using, for operating and making use of, whether
      personally or on behalf of a third party, remunerated or free of
      charge, the Contributions herein assigned, for the purposes of
      carrying out any form of processing, to whatever end.
    \item It is agreed that all of the rights considered above are
      assigned for the duration of copyright, for the whole world.
  \end{itemize}
  It is specified that any author-employee of the Contributor remains
  the holder of his or her moral right over said Contributions.
  \item In the event that the Contributor wishes to cease submitting
    further Contributions, it promises to inform GNUnet e.V. of this
    in advance by registered mail with return receipt. It further
    promises to ensure that all participation by its employees in the
    development of the Software ceases.
  \item The Contributor and its employees nevertheless remain free to
    use the Software under the conditions provided for by the GNU AGPL
    v3 license or any other license which may subsequently become
    attached to it.
  \item The assignment is non-exclusive. In particular, the
    Contributor remains free to make the same Contributions to other
    projects and under other licenses.  However, this exception only
    applies to the Contribution, and in particular does not relieve
    the Contributor from the obligations of the GNU AGPLv3+ or other
    licenses of the Software, in cases where these licenses may
    otherwise apply to the Contribution.
  \item GNUnet e.V. remains free to terminate the submission of
    Contributions by the Contributor at any time.
\end{enumpar}

\section{GUARANTEES}

\begin{enumpar}
 \item The Contributor guarantees that it possesses copyrights over
   the Contributions (non-exclusively) transferred and that it holds
   all of the rights necessary to assign said Contributions to GNUnet
   e.V.. The Contributor guarantees, in particular, that the
   Contributions assigned to GNUnet e.V.  do not infringe third
   parties' intellectual property rights, including those protected by
   patent law.  Should the Contributor become aware, after the
   assignment of a Contribution to GNUnet e.V., of details calling
   into question this article and the validity of the assignment
   covered by this contract, the Contributor promises to inform GNUnet
   e.V. of this promptly, in order that the latter may withdraw the
   Contribution or Contributions at issue from the Software, as
   necessary.
 \item The Contributor acknowledges that it holds no patent which
   could be enforced against any use by GNUnet e.V. of the copyright over
   the Contributions which are assigned to it. In the opposite case,
   the Contributor promises to raise no objection to GNUnet e.V., or any of
   its licensees, sub- licensees or assignees, on the grounds of the
   use of its patent in the Contributions assigned.
 \item The Contributor is bound by no obligation to maintain its
   Contributions, except as voluntarily undertaken on its
   part. Moreover, the Contributor declares that it is assigning its
   Contributions ``as they are'', without guarantee as to their
   commercial value, and without guaranteeing GNUnet e.V. that the
   Contributions are free of errors or correspond to its needs.
 \item GNUnet e.V. guarantees that it will continue to distribute
   all future versions of the Software under a free software license,
   as defined by the Free Software Foundation (\url{http://www.fsf.org/}).
\end{enumpar}

\section{MISCELLANEOUS}

\begin{enumpar}
 \item This contract is subject to Swiss law.
 \item Any dispute concerning the interpretation, validity or
   execution of this contract will be submitted, failing an
   out-of-court resolution, to the competent German court.
 \item If one or several stipulations of this contract are held to be
   invalid or declared such in application of a law or regulation, or
   by reason of a final ruling by a competent court, the others will
   retain all of their force and scope.
\end{enumpar}



\end{document}
